% Proved after complete root and companion audits; exact candidate retained.
\subsection{Two shared cyclic quotients force the full first layer at \(p=5\)}

\begin{lemma}[A coordinate whose high-displacement line is not the last line]
\label{lem:B-RD87-coefficient}
\status{Proved; full independent root and companion audits passed.}
Let \(E/L\) be a totally ramified Galois extension of
complete discretely valued mixed-characteristic fields
with group \(G\simeq\mathbf F_p^2\) and the same perfect
residue field. Suppose \(G\) has two lower breaks
\(t<u\), with last ramification line \(T\).
Let \(x\in\mathfrak m_E\), put \(A=v_E(x)\), and
suppose \(v_p(A)=1\). Assume a line \(H\subset G\)
and integers \(D_0<D_1<v_E(p)\) satisfy
\[
 v_E(gx-x)=
 \begin{cases}
  D_1,&0\ne g\in H,\\
  D_0,&g\notin H.
 \end{cases}                                         \tag{B-RD87.1}
\]
If \(H\ne T\), then
\[
     \frac{u+(p-1)t}{p}=\frac{D_0-A}{p}.              \tag{B-RD87.2}
\]
The left side is the lower break of any cyclic
degree-\(p\) quotient in which \(T\) has nonzero image.
\end{lemma}
\begin{proof}
Choose a uniformizer \(\varpi\) of \(E\) and fixed
residue representatives in \(\mathcal O_L\). Write
\(x=\sum_{j\ge A}a_j\varpi^j\), with \(a_A\) a unit
and all nonzero \(a_j\) units. Let \(J\) be the first
degree with \(p\nmid J\) and \(a_J\ne0\), allowing
\(J=+\infty\).

Up to valuation \(D_1<v_E(p)\), terms whose integer
binomial coefficient is divisible by \(p\) do not
contribute. If an automorphism has lower break \(b\),
the leading degree \(A\), whose \(p\)-adic order is
exactly one, first changes in degree \(A+pb\), with
nonzero coefficient. Every higher degree divisible
by \(p\) changes later. The first prime-to-\(p\)
degree first changes in degree \(J+b\). Therefore
the first possible change has degree
\[
                 m(b)=\min\{A+pb,J+b\},              \tag{B-RD87.3}
\]
and it is the actual change degree when the two
quantities are distinct and below \(v_E(p)\).
Only equality of the two quantities permits
cancellation. These assertions follow directly
by expanding \(g\varpi=\varpi(1+d\varpi^b+\cdots)\);
the same fixed-coefficient convention is used in B60
and B81.

Take \(0\ne g\in H\) and \(0\ne h\in T\).
Since the lines differ, their breaks are \(t,u\),
and their coordinate displacements are \(D_1,D_0\).
In particular \(m(u)\le D_0\). Both expressions in
(B-RD87.3) strictly increase with \(b\), so
\(m(t)<D_0<D_1\). The first change for \(g\) must
cancel. Consequently \(J\) is finite and
\[
                      J=A+(p-1)t.                    \tag{B-RD87.4}
\]
At the larger break \(u\), the degree \(J+u\) is
strictly less than \(A+pu\), hence is the unique
first change. Thus
\(D_0=J+u=A+(p-1)t+u\), proving the identity.

For the quotient assertion, every nonzero coset of
a kernel line different from \(T\) contains one
element of \(T\) and \(p-1\) elements outside it.
The exact norm-polynomial quotient formula
\cite[Proposition 3.7 and full proof,
pp.~14--15]{ColmezCorpsLocaux} averages their
\(b+1\) values and gives precisely
\((u+(p-1)t)/p\). The quotient killing \(T\) instead
has break \(t\). Thus the former is the larger
of these two cyclic quotient breaks.
\end{proof}

Now let \(K\) be complete discretely valued of mixed
characteristic \((0,5)\), with algebraically closed
residue field, \(e_K=20\), and \(\mu_5\subset K\).
Let \(L/K\) be the common tame good field of degree
24. Suppose two whole potentially supersingular
height-two groups have isomorphic actual
\(125\)-torsion, distinct finite special characters,
and tame reducible residual inertia.

\begin{lemma}[The two point fields share a diagonal cyclic quotient]
\label{lem:B-RD87-fields}
\status{Proved; full independent root and companion audits passed.}
If \(\dim W_1=3\), then there are two actual Galois
point fields \(E_1,E_2\) of degree 25 over \(L\),
sharing a diagonal cyclic degree-five quotient field
\(D\). Write \(d\) for the break of \(D/L\), and
\(o_i\) for that of the complementary off-diagonal
cyclic quotient of \(E_i/L\). Then
\[
     d>0,\quad d\equiv0\pmod{24},\qquad
        o_1\equiv6\pmod{24},\quad o_2\equiv18\pmod{24}.
                                                        \tag{B-RD87.5}
\]
Each \(E_i/L\) has two distinct lower breaks, with
last ramification line diagonal if \(d>o_i\), and
off-diagonal if \(o_i>d\).
\end{lemma}
\begin{proof}
By B84, after interchanging the groups and making
the justified common finite twist, the parameters
are \(B_1=25\), \(B_2=75\). Put the high lines at
\(H_1=\mathbf F_5e_1\), \(H_2=\mathbf F_5e_2\).
The residual tame eigenvalue exponents on these
two vectors are \((0,6)\), or \((12,18)\),
modulo 24. Their ratio is the same in both rows.
Thus tame conjugation has weights 6 on \(E_{21}\),
18 on \(E_{12}\), and zero on diagonal matrices.

B83 and B85 give
\[
 W_1=\mathbf F_5E_{12}\oplus\mathbf F_5E_{21}
                    \oplus\mathbf F_5Z,
             \qquad Z=\operatorname{diag}(a,b),\quad ab\ne0.
                                                        \tag{B-RD87.6}
\]
Indeed the tame spectral projectors separate all
three weights, and both diagonal-coordinate
projections are nonzero. The common second-level
division field over \(L\) has additive Galois group
exactly \(W_1\).

Choose a primitive second-level point above \(e_i\),
and call its field \(E_i\). First-level rationality
makes it Galois, with translation group \(W_1e_i\).
Since \(a,b\ne0\), both evaluations are surjective,
with kernels \(\mathbf F_5E_{12}\) for \(i=1\)
and \(\mathbf F_5E_{21}\) for \(i=2\).
Their common quotient by both off-diagonal axes
is the same degree-five field \(D\). In each
translation group the diagonal line is the image
of \(Z\), and the other line is the image of the
remaining off-diagonal matrix. All these kernel
spaces are stable under tame conjugation, so the
three cyclic quotient fields are Galois over \(K\)
as well as over \(L\).

We spell out the tame-break congruence. In a cyclic
degree-five quotient \(F/L\), let \(\tau\) lift
the tame generator of \(L/K\), whose leading
uniformizer multiplier is a primitive 24th root
\(\zeta\). On a uniformizer of \(F\), its leading
multiplier is \(\eta\), with \(\eta^5=\zeta\),
so \(\eta=\zeta^5\). If the nonidentity wild
elements have lower break \(b_F\), conjugating
\(g\varpi=\varpi+c\varpi^{b_F+1}+\cdots\)
shows that \(\tau\) acts on their leading
coefficients by \(\eta^{b_F}=\zeta^{5b_F}\).
The tame weights of the three quotient groups
are respectively \(0,6,18\). This gives
(B-RD87.5), since multiplication by five fixes
6 and 18 modulo 24.

Ramification groups are characteristic under
field automorphisms and hence tame-stable.
The two distinct tame characters in each
\(\operatorname{Gal}(E_i/L)\) imply that a proper
ramification line is one of these two axes.
There cannot be just one break: then the
norm-polynomial average would give equal breaks
for all cyclic quotients, contrary to (B-RD87.5).
For two breaks \(t<u\), the quotient killing the
last line has break \(t\), and the other displayed
quotient has the strictly larger break
\((u+4t)/5\). This proves the final assertion.
\end{proof}

\begin{proposition}[Three first-layer dimensions are impossible at \(p=5\)]
\label{prop:B-RD87-three}
\status{Proved; full independent root and companion audits passed.}
Under the stated actual \(125\)-comparison hypotheses,
one has \(\dim W_1\ne3\).
\end{proposition}
\begin{proof}
Suppose otherwise and apply the preceding lemma.
In the two good models the first-level low/high
valuations over \(L\) are \((5,95)\) and
\((15,45)\), respectively. Each parameter is
greater than \(e_K=20\), so all inverse roots over
a high target have valuation equal to that
target's valuation divided by 25, as in B85.
The same single-segment assertion holds over a
low target. Since \([E_i:L]=25\), their actual
primitive coordinates have the following
integral valuations and translation displacements:
\[
\begin{array}{c|c|r|r|r|c}
 \text{field}&\text{model}&A&D_0&D_1&H\\\hline
 E_1&1&95&125&2375&\text{diagonal}\\
 E_1&2&15&375&1125&\text{off-diagonal}\\
 E_2&1&5&125&2375&\text{off-diagonal}\\
 E_2&2&45&375&1125&\text{diagonal}
\end{array}                                             \tag{B-RD87.7}
\]
Here \(A\) is the point-coordinate valuation;
\(D_0,D_1\) are the low/high values of
\(v_{E_i}(gT(P)-T(P))\); and \(H\) is the line
of translations attaining \(D_1\). These are
actual point fields, shared under the given
torsion isomorphism. Formal addition multiplies
the translating coordinate by a unit, so the
displacement entries are exactly 25 times the
first-level valuations. Every \(A\) has
5-adic order one, and
\[
           v_{E_i}(5)=25\cdot24\cdot20=12000>2375.
                                                        \tag{B-RD87.8}
\]
Thus Lemma~\ref{lem:B-RD87-coefficient} applies to
each row without a mixed-characteristic equality
or truncation exception.

In \(E_1\), if the last ramification line were
off-diagonal, the first row would force its
larger cyclic quotient break to be
\(o_1=(125-95)/5=6\). But it would have to
exceed \(d\ge24\), a contradiction. Hence the
last line in \(E_1\) is diagonal. The second
row now gives
\[
                         d=(375-15)/5=72.             \tag{B-RD87.9}
\]
In \(E_2\), a diagonal last line would make
the third row give \(d=(125-5)/5=24\), contrary
to (B-RD87.9). An off-diagonal last line would
make the fourth row give
\(o_2=(375-45)/5=66\), which would have to
exceed \(d=72\), again impossible. These are
the only two ramification-line choices. The
assumed three-dimensional first layer therefore
cannot occur.
\end{proof}

\begin{corollary}[The tame residual ambiguity requires the full congruence kernel]
\label{cor:B-RD87-full}
\status{Proved; full independent root and companion audits passed.}
At \(p=5,e_K=20\), a different-special-character
actual \(125\)-congruent pair in the tame reducible
residual sector must have
\[
    W_1=W_2=M_2(\mathbf F_5),\qquad
    W=1+5M_2(\mathbf Z/125),\qquad |W|=5^8.           \tag{B-RD87.10}
\]
Its base field does not contain \(\mu_{25}\).
The necessary parameter pattern remains
\(\{25,75\}\), type A/B up to common quadratic twist.
\end{corollary}
\begin{proof}
B85 gives \(\dim W_1\ge3\), and the proposition
excludes equality. The ambient matrix space has
dimension four, so \(W_1=M_2(\mathbf F_5)\).
Actual fifth powers give \(W_1\subset W_2\),
so the terminal layer is the same full space.
Thus \(W\) has order \(5^8\), equal to that of
the full principal congruence kernel containing
it, and the groups are equal. The determinant
modulo 25 has nonzero trace on the first layer.
On \(L\) this is the actual cyclotomic character,
so \(\mu_{25}\) is not contained in \(L\), and
in particular not in \(K\). The parameter and
finite-type statement is precisely B84.
\end{proof}

\begin{remark}[Scope and the shared-quotient requirement]
\label{rem:B-RD87-scope}
\status{Proved; full independent root and companion audits passed.}
The statement concerns only the tame reducible
residual sector of the \(p=5,e_K=20\) problem.
It does not bound the other residual sectors,
construct the surviving full congruence image,
or settle the unrestricted least index, which
remains in \(\{2,3\}\).
The shared diagonal field \(D\) follows from
the actual same three-dimensional first layer.
Two arbitrary degree-25 extensions with
unrelated diagonal quotients would not give
this contradiction. No cyclotomic identification
of \(D\) is assumed.

The coefficient proof keeps the actual point
lifts fixed, retains the strict bound
\(D_1<v_{E_i}(5)\), and includes the possibility
that the first prime-to-five coefficient is
absent until cancellation forces it to exist.
The complete B60 and B81 fixed-coefficient
proofs and the B83--B85 actual field and
matrix arguments were reread. Serre's original
Section 3.1, pp.~126--128, was reopened and
read for the ramification definitions and
filtration conventions; its quotient theorem
refers to other proofs. The complete direct
norm-polynomial proof used here was freshly
read in Colmez, Proposition 3.7, pp.~14--15.
The tame leading-coefficient calculation is
written explicitly above. No unexamined
finite-residue assertion is imported.

The retained exact supplement is
\path{code/reduction_descent/verify_p5_tame_full_first_layer.py}.
The author run passed at 01:32:46.564572 UTC on
20 November 2026 in 0.03949103799823206 seconds:
three symbolic identities, 900 literal binomial
checks, eight strict inverse-polygon vertices,
all four coordinate rows, all 2,000 vectors in
the 16 possible diagonal first-layer spaces,
three exact tame congruences, 12 invariant-line
checks, and all four last-line alternatives.
The uniform coefficient argument and the actual
field constructions are proved above rather
than extrapolated from this finite supplement.
No failed mathematical computation preceded
this run; the exact executed code and output
are retained without replacement.

\smallskip
\noindent\emph{Independent audits.}
The complete root and companion proof, source and code audits passed.
Their literal author-output-preserving reruns took
\(0.037566576997051015\) and
\(0.034868966991780326\) seconds, respectively.
The exact audited candidate, per-auditor read scopes,
all code and outputs, and the promotion diff are retained.
\end{remark}
